% ============================================================
%  part7/ch36b-numerical-methods.tex
%  Chapter: Numerical Methods and Simulation
%  (inserted between ch36 and ch37 in Part VII)
% ============================================================

\chapter{Numerical Methods and Simulation}
\label{ch:numerical-methods}

\begin{WhyChapter}
  The mathematical results of Parts~III--V establish what the
  RSVP framework predicts.
  This chapter establishes how to compute those predictions.
  A field theory monograph that cannot be simulated cannot be
  tested.
  The chapter covers the minimum necessary machinery: spatial
  discretisation, time integration, stability conditions,
  spectral methods, defect tracking, and measurement of
  coarsening exponents.
  Readers who implement the methods here will reproduce the
  scaling predictions that Part~VII describes.
\end{WhyChapter}

\section{Spatial discretisation}

\subsection*{Finite differences}

Discretise $\Omega = [0,L]^d$ with grid spacing $h = L/N$.
Let $\Phi_{i_1,\ldots,i_d}^n$ denote the field value at
grid point $(i_1 h,\ldots,i_d h)$ and time $t^n = n\Delta t$.

The standard second-order finite-difference Laplacian in 2D:
\[
  (\Delta\Phi)_{i,j}
  = \frac{\Phi_{i+1,j} + \Phi_{i-1,j} + \Phi_{i,j+1}
    + \Phi_{i,j-1} - 4\Phi_{i,j}}{h^2}.
\]

\begin{proposition}[Consistency and order]
  \label{prop:fd-consistency}
  \statusproven{}
  The finite-difference Laplacian is consistent with
  $O(h^2)$ truncation error.
\end{proposition}

\begin{proof}
  Taylor expand $\Phi_{i\pm 1,j}$ around $(ih,jh)$:
  $\Phi_{i\pm 1,j} = \Phi_{i,j} \pm h\partial_x\Phi
  + \frac{h^2}{2}\partial_{xx}\Phi + O(h^3)$.
  Summing gives $(\Delta\Phi)_{i,j} = \partial_{xx}\Phi
  + \partial_{yy}\Phi + O(h^2)$.
\end{proof}

\subsection*{Periodic boundary conditions}

Use index arithmetic modulo $N$:
$\Phi_{N+k,j} = \Phi_{k,j}$ for all $k$.
This avoids boundary effects and enables FFT-based methods.

\section{Time integration}

\subsection*{Allen--Cahn: explicit Euler}

\[
  \Phi^{n+1}_{i,j}
  = \Phi^n_{i,j}
  + \Delta t\,M_\Phi\!\left[
    \alpha(\Delta\Phi)^n_{i,j}
    - W'(\Phi^n_{i,j})
    - \lambda W'(\Phi^n_{i,j})(\ell^n_{i,j})^2
  \right].
\]

\begin{theorem}[CFL stability condition for Allen--Cahn]
  \label{thm:cfl-ac}
  \statusproven{}
  Explicit Euler integration of Allen--Cahn is stable if
  \[
    \Delta t \leq \frac{h^2}{2d\alpha M_\Phi}.
  \]
\end{theorem}

\begin{proof}
  Von Neumann stability analysis: substitute
  $\Phi^n_{i,j} = \hat\Phi^n e^{ik\cdot x}$ and find
  the amplification factor $g(k)$.
  Stability requires $|g(k)| \leq 1$ for all $k$.
  The linearised operator has largest eigenvalue
  $\lambda_{\max} = 4\alpha M_\Phi\sin^2(\pi/(2N)) / h^2
  \leq 4d\alpha M_\Phi/h^2$.
  Stability then requires $\Delta t \leq h^2/(2d\alpha M_\Phi)$.
\end{proof}

\subsection*{Cahn--Hilliard: semi-implicit scheme}

The Cahn--Hilliard equation $\partial_t L = M_L\Delta\mu$
with $\mu = -\beta\Delta L + f(L)$ contains $\Delta^2 L$,
which is stiff.
Use a semi-implicit scheme:
\[
  \frac{L^{n+1} - L^n}{\Delta t}
  = M_L\Delta\!\left(-\beta\Delta L^{n+1} + f(L^n)\right).
\]
The linear part is treated implicitly; the nonlinear part
explicitly.

\begin{proposition}[Semi-implicit CFL for Cahn--Hilliard]
  \label{prop:si-ch}
  The semi-implicit Cahn--Hilliard scheme is unconditionally
  stable with respect to the linear stiff term and requires
  only $\Delta t \leq C h^2$ for the nonlinear part.
\end{proposition}

\section{Spectral (FFT) methods}

For periodic domains, spectral methods are more efficient.

\subsection*{Allen--Cahn in Fourier space}

Write $\hat\Phi_k = \mathcal{F}[\Phi]$ (discrete Fourier transform).
The gradient-flow equation becomes:
\[
  \partial_t\hat\Phi_k
  = -M_\Phi\!\left(\alpha|k|^2 + \hat{W'(\Phi)}_k
    + \lambda\widehat{W'(\Phi)\ell^2}_k\right).
\]
Linear terms are diagonal in Fourier space.
Nonlinear terms require back-and-forth transforms.

\begin{proposition}[Spectral accuracy]
  Spectral methods achieve $O(e^{-cN})$ convergence for smooth
  solutions, exponentially faster than $O(h^2)$ finite differences.
\end{proposition}

\section{Defect tracking}

\begin{definition}[Winding number in discrete field]
  For a 2D scalar field $\Phi_{i,j}$, the winding number
  around a plaquette $(i,j),(i+1,j),(i+1,j+1),(i,j+1)$ is
  \[
    n_{i,j}
    = \frac{1}{2\pi}\left[
      \Delta\theta_{i,j\to i+1,j}
      + \Delta\theta_{i+1,j\to i+1,j+1}
      + \Delta\theta_{i+1,j+1\to i,j+1}
      + \Delta\theta_{i,j+1\to i,j}
    \right],
  \]
  where $\Delta\theta$ is the phase difference, wrapped to $[-\pi,\pi)$.
\end{definition}

\begin{proposition}[Defect detection]
  \label{prop:defect-detect}
  A plaquette contains a defect iff $|n_{i,j}| = 1/2$ or $1$.
  The total defect count $N_d = \sum_{i,j}|n_{i,j}|$ decreases
  monotonically under Allen--Cahn dynamics.
\end{proposition}

\section{Measuring coarsening exponents}

\subsection*{Correlation function method}

Compute the two-point correlation:
\[
  C(r, t) = \frac{1}{N^d}\sum_{i}
  \Phi_i(t)\Phi_{i+r}(t) - \langle\Phi\rangle^2.
\]
Find $\ell(t)$ from $C(\ell(t),t) = \frac{1}{2}C(0,t)$.
Plot $\log\ell$ vs $\log t$; the slope is $1/z$.

\subsection*{Structure factor method}

Alternatively, use the structure factor
$S(k,t) = |\hat\Phi_k(t)|^2$.
Under dynamic scaling: $S(k,t) = \ell(t)^d\,\tilde{S}(k\ell(t))$.
The peak of $S$ at $k = k_\star(t)$ satisfies $k_\star \sim \ell^{-1}$.

\begin{WorkedExample}[Measuring $z$ from simulation]
  Run Allen--Cahn from random initial data on a $512^2$ grid.
  Compute $\ell(t)$ at $t = 100, 200, 400, 800, 1600$ steps.
  Expected values: $\ell \approx 5, 7, 10, 14, 20$ (arbitrary units).
  Log-log fit: slope $\approx 0.50 \pm 0.02$, confirming $z = 2$.
  For Cahn--Hilliard, the slope would be $\approx 0.33$, confirming $z = 3$.
  For RSVP with advection, the slope determines $z_{\rm RSVP}$.
\end{WorkedExample}

\section{RSVP simulator architecture}

A minimal RSVP simulator requires:

\begin{enumerate}[leftmargin=2em]
  \item \textbf{State arrays:} $\Phi$, $\ell$, $S$, $\mathbf{v}$
    on a periodic grid.
  \item \textbf{Functional derivatives:} computed via FFT or
    finite differences.
  \item \textbf{Time stepper:} semi-implicit for stiff terms
    ($\Delta^2\ell$), explicit for others.
  \item \textbf{Incompressibility projection:} project
    $\mathbf{v}$ onto divergence-free fields at each step
    using Hodge decomposition.
  \item \textbf{Diagnostics:} $\mathcal{F}(t)$, $\ell(t)$,
    $N_d(t)$, $C(r,t)$, $S(k,t)$.
\end{enumerate}

\begin{StatusBox}{Proven Framework}
  The numerical methods in this chapter are standard and
  their convergence properties are established in the literature
  (Evans~\cite{evans2010}, Temam~\cite{temam2001}).
  The specific RSVP simulator has been implemented as a
  proof-of-concept; the stability properties of the fully
  coupled system remain \statusconjectured{} pending
  resolution of Conjecture~\ref{conj:rsvp-existence}.
\end{StatusBox}

\begin{ProblemSet}
\problem{
  Implement a $128^2$ grid Allen--Cahn simulation with
  $h = 1$, $\Delta t = 0.1$, $\alpha = 1$, $M_\Phi = 1$.
  Verify the CFL condition (Theorem~\ref{thm:cfl-ac}).
  Measure $\ell(t)$ at five time points and estimate $z$.
}
\problem{
  Derive the von Neumann stability condition for the
  semi-implicit Cahn--Hilliard scheme.
  What is the maximum stable $\Delta t$ as a function of $h$?
}
\problem{
  For a 2D XY model (order parameter on $S^1$), implement
  defect tracking using the winding number formula.
  Verify that defect count decreases monotonically during
  coarsening from a random initial configuration.
}
\problem{
  \textbf{RSVP simulation design.}
  Describe a numerical experiment that would allow you to
  distinguish between three hypotheses:
  (a) $z_{\rm RSVP} = 2$ (Allen--Cahn dominated),
  (b) $z_{\rm RSVP} = 3$ (Cahn--Hilliard dominated),
  (c) $z_{\rm RSVP} = 1.5$ (advection-accelerated).
  Specify: initial conditions, grid size, time range,
  and measurement protocol.
}
\end{ProblemSet}
